% !TeX spellcheck = en_us
\documentclass[a4paper,12pt,fleqn]{article}

\usepackage[left=1.50cm, right=1.50cm, top=3cm, bottom=2cm, 
headheight= 2.5cm, headsep= 4pt, footskip=25pt]{geometry}
\usepackage{fancyhdr}
\usepackage{lmodern}
\usepackage[table]{xcolor}
\usepackage{amssymb,amstext,amsmath}
\usepackage{booktabs}
\usepackage{bm}
\usepackage{setspace}
\usepackage{graphicx}
\usepackage[colorlinks]{hyperref}
\setlength{\parindent}{0mm}
\setlength{\parskip}{2mm}
\usepackage{enumitem}
\usepackage{titlesec}
\usepackage{gensymb}
\usepackage{pgfgantt}
\usepackage{lscape}
\usepackage[absolute]{textpos}
\usepackage{comment}
\usepackage{cite}
\usepackage{tikz}
\usepackage{pgfplots}
\usetikzlibrary{shapes.geometric}
\usepackage{array}
\newcolumntype{P}[1]{>{\centering\arraybackslash}p{#1}}
\usepackage{longtable}
\usepackage[font={small}]{caption}
\usepackage[most]{tcolorbox}
\usepackage{tikz}
\usepackage{stmaryrd}
\usepackage{wrapfig}
%\usepackage[none]{hyphenat}
\usepackage{ctable}
\usepackage{ragged2e}
\usepackage[official]{eurosym}
\newcommand*\circled[1]{\tikz[baseline=(char.base)]{
		\node[shape=circle,draw,inner sep=1pt] (char) {#1};}}
\usepackage[symbol]{footmisc}
\renewcommand{\thefootnote}{\fnsymbol{footnote}}
\renewcommand{\thempfootnote}{\fnsymbol{mpfootnote}}
\renewcommand*\footnoterule{}

\usepackage{xcolor}
\definecolor{docColor}{cmyk}{1.00,1.00,0, 0.40}  % 0.90\% of black
\color{docColor}

\newtcolorbox{myhlbox}[2][]{colback=docColor!05, colframe=docColor!20, width=1.0\linewidth,
	boxrule=0pt, left=1em, right=1em, top=1em, bottom=1em,
	colbacktitle=docColor!20, coltitle=docColor,enhanced, attach boxed title to top left={xshift=2mm, yshift=-1mm},
	title=#2, #1}

\hypersetup{
	colorlinks=true,
	citecolor=blue,
	filecolor=magenta,
	linkcolor=blue,
	urlcolor=blue,
	pdftitle={SeminarAnnouncements},
	%bookmarks=true,,
}

\graphicspath{{./figs/}}

\definecolor{dgreen}{rgb}{0.0, 0.5, 0.0}

\allowdisplaybreaks

\newcommand*\mystrut[1]{\vrule width0pt height0pt depth#1\relax}
\renewcommand{\thesection}{\arabic{section}.}
\renewcommand{\thesubsection}{\arabic{section}.\arabic{subsection}}
\renewcommand{\thesubsubsection}{\arabic{section}.\arabic{subsection}.\arabic{subsubsection}}
\titleformat{\section}{\sffamily\large\bfseries}{\thesection}{1.0em}{}
\titleformat{\subsection}{\sffamily\large\bfseries}{\thesubsection}{1.0em}{}
\titleformat{\subsubsection}{\sffamily\normalsize\bfseries}{\thesubsubsection}{1.0em}{}
\titlespacing\section{0pt}{0pt plus 2pt minus 2pt}{0pt plus 2pt minus 2pt}
\renewcommand{\baselinestretch}{1.00}

\newcommand\CircledImage[1]{%
	\tikz\path[fill overzoom image={#1}, line width=0mm]circle[radius=0.5\linewidth];%
}

\input{format/defs}

\fancypagestyle{titlepage}{
	\fancyhf{}
	\insertLXheader
}

\fancypagestyle{plain}{%
	\fancyhf{}
	\renewcommand{\headrulewidth}{0pt}
	\renewcommand{\footrulewidth}{0pt}
	\insertnormalpage
}

\pagestyle{plain}

\begin{document}
	\sffamily
\thispagestyle{titlepage}
\vspace*{-2em}
\begin{center}
	\huge \textbf{Mechanics Seminar series 2025 -- 26}
\end{center}

\begin{center}
	\Large 
	\textbf{Modeling the Unique Thermomechanical Behaviors of Liquid Crystal Elastomers} 
\end{center}
\begin{center}
	\large
	Thao (Vicky) Nguyen \\ 
	{\large Department of Mechanical Engineering \\
		Johns Hopkins University
	} \\
	\vspace*{1em}
	\textbf{Date, time, and venue }\\
	January 08, 2026 (2 -- 3 pm), Amphithéâtre Pole Mecanique
\end{center}

\begin{myhlbox}{{\large Abstract}}
	Liquid crystal elastomers (LCEs) exhibit complex thermomechanical behaviors that can be harnessed for a wide range of applications in soft robotics, biomedical devices, and energy absorption. The material consists of stiff mesogens bound within an elastomeric network of flexible polymer chains. The mesogens interact energetically and can order and disorder in response to temperature and mechanical deformation, amongst other stimuli. This enables LCEs to undergo reversible phase transitions between the disordered isotropic, ordered monodomain, and polydomain states. The motion of the mesogens relative to the polymer network results in unique behaviors, including reversible actuation response to temperature and soft elasticity. Additionally, LCEs display enhanced energy dissipation compared to conventional elastomers due to the viscous rotation of the mesogens and the relaxation of the network chains. These viscoelastic dissipation mechanisms can be utilized to design LCE materials and structures with extraordinary toughness, impact energy absorption, and mechanical damping. However, these same properties may hinder the actuation and morphing capabilities of the material.  Predictive modeling is essential for efficiently designing and optimizing LCE structures to attain the desired performance. In this presentation, I will outline our recent efforts to develop generalized continuum theories for the thermomechanical behavior of monodomain nematic elastomers that incorporate the rate-dependent deformation mechanisms of the mesogens and the network chains. I will demonstrate the predictive capabilities of the theories and highlight their application to design energy-absorbing architected materials, assess the effectiveness of actuators, and optimize the director pattern of LCE structures to maximize viscoelastic dissipation.   
\end{myhlbox}

\vspace*{1em}

\begin{myhlbox}{{\large About the speaker}}
	\begin{wrapfigure}{r}{0.25\textwidth}
		\vspace*{-2em}
		\CircledImage{images/vicky}
	\end{wrapfigure}

	Thao (Vicky) Nguyen received her S.B. from MIT in 1998, and M.S. and Ph.D. from Stanford in 2004, all in mechanical engineering. She was a research scientist at Sandia National Laboratories in Livermore from 2004- 2007 before joining Johns Hopkins University, where she is currently a Professor in Mechanical Engineering with secondary appointments in Materials Science and Engineering and Ophthalmology. She is also Deputy Director of the Hopkins Extreme Materials Institute. Dr. Nguyen’s research encompasses the mechanics of soft tissues, stimuli-responsive soft materials, and engineering polymers. Her lab currently studies the biomechanics of the optic nerve head in glaucoma, the mechanics of recycled polymers, and the mechanical behavior of liquid crystal elastomers and DNA hydrogels. Dr. Nguyen has received numerous awards, including the  2008 Presidential Early Career Award for Scientists and Engineers (PECASE) and the NNSA Office of Defense Programs Early Career Scientists and Engineers Awards for her work on modeling the thermomechanical behavior of shape memory polymers. In 2013, she received the NSF CAREER award for studying the micromechanisms of growth and remodeling of collagenous tissues, the Eshelby Mechanics Award for Young Faculty, and the Sia Nemat-Nasser Early Career Medal from the Materials Division of ASME.  She received the T.J.R. Hughes Young Investigator Award from the ASME Applied Mechanics Division in 2015, the Van C. Mow Medal from the ASME Bioengineering Division in 2024, the James R. Rice Medal from the Society of Engineering Science in 2025, and the Centennial Mid-Career Award from the Materials Division of ASME also in 2025.  She is a Fellow of ASME and the American Institute for Medical and Biological Engineering (AIMBE). She served as the President of the Society of Engineering Science (SES) in 2020 and is currently the  Editor-in-Chief of the Journal of Biomechanical Engineering.  

\end{myhlbox}
\vfill

\end{document}